\section{问题分析与总体建模思路}

\subsection{赛题多维度系统特征剖析}

本赛题涵盖三维地理地形建模、非线性飞行动力学、带时间窗车辆路径优化（VRPTW）、无线电电磁传播及应急资源配置等多个领域。分析各子模块的物理机制与时空约束是构建模型的基础：

\begin{enumerate}
    \item \textbf{复杂三维地形与净空飞行剖面分析}：
    受灾区域位于广西横州镇龙乡，属于典型的喀斯特地貌与山地丘陵过渡带。官方提供的 $1486 \times 1309$ 栅格 30 米高精度 DEM 矩阵表明，全域地面高程在 50 米至 750 米之间剧烈波动。无人机在任意两个地理节点间飞行时，必须保证高于沿途地貌最高点的安全净空高度（取 50 米安全裕度）。因此，航段飞行不是平直欧氏空间位移，而必须精确分解为垂直爬升至巡航高度、定高水平巡航以及下降至目标投放点的完整剖面\cite{zhang2021uav}。高起伏地形直接导致正反向航段的爬升与下降高度产生显著不对称性。
    
    \item \textbf{非线性动力学能耗与最大载荷上界分析}：
    多旋翼无人机的升力产生依赖旋翼高速旋转向下排气，随着有效载荷 $w$ 的增加，机身总重增大导致所需诱导功率急剧上升。官方附录给出的 $(3/2)$ 次方等效航程递减公式表明，单位距离的放电损耗关于载重具有超线性增长特征。同时，爬升重力势能做功以逆变器与电机效率 $\eta_{climb} = 0.72$ 计入能耗，而下降过程不回收重力势能。因此，在固定的返航安全剩余电量余量 $\eta$ 下，必须精确搜索各机型针对各服务区的物理有效载荷上界，防止因过载引发航中失速坠机。
    
    \item \textbf{异构机型与车电分离双时间轴耦合分析}：
    题目提供了 A、B、C 三类性能各异的多旋翼运输机以及专用的锂电池组。其中，机身实体（4 架 A、2 架 B、2 架 C）与共享电池（6 组 A、4 组 B、4 组 C）在物理上相互解耦。更关键的是，大容量动力电池的恒流-恒压充电遵循典型的两阶段非线性规律（SOC 达到 0.90 之前耗费 65\% 时间，0.90 至 1.00 阶段耗费 35\% 时间）。这意味着实体机返航后可通过更换已充饱的待命电池迅速实施下一架次作业，而放电后的电池则需进入独立的充电排队流水线。调度决策不仅要对 8 架飞机进行排程，还必须严格维持 14 组电池的时间轴防冲突，构成极具挑战性的多实体并发排程问题。
    
    \item \textbf{电磁视距遮挡与空中立体中继拓扑分析}：
    无线电高频信号（$f=2400\text{ MHz}$）具有典型的视距传播特征，在遭遇山脊实体刺入视线时会产生至少 10 dB 的强烈阴影衰落。在复杂的山地沟壑中，大量低海拔受灾点（如深谷中的 S003、S007 等）与位于 O01 的固定网关 G01 之间存在严重的视线遮挡。为达成“起飞至降落全阶段 100\% 连续通信”的严苛指标，必须在山脊制高点部署长续航中继无人机。中继机不仅要维持高空悬停与通信载荷供电，还受到 300 米离地高度限制与地面换电维护工时约束，需与运输机群形成高精度的时空编队协同。
    
    \item \textbf{航线强绑定与管理分区连通分支机理分析}：
    在第四问的应急责任分区中，题目给出了关键性的业务约束：“若同一运输航次涉及多个服务区，则这些服务区必须划入同一任务组”。这一刚性规则使得多点回路航线在受灾点之间建立了等价相依关系。在代数图论中，由航线边构成的服务区关联图必然能够分解为若干相互独立的连通分支（Connected Components）。各任务组封闭独立执行切断了机型与电池的跨区时分复用，从而诱发了装备峰值并发需求的激增与刚性短缺。
\end{enumerate}

\subsection{总体建模思路与技术路线图}

针对上述多维度系统特征，本文确立了“三维解析为基、分阶递进优化、车电异步协同、图论拓扑定界”的总体建模思路。

首先，基于高精度地理数据全面构建灾区地理信息与网络拓扑。如图~\ref{fig:terrain}所示，本文基于 30m DEM 矩阵构建高精度矢量等高线测绘底图，精准刻画横州镇龙乡高山深谷的地表起伏特征，并标定了调度中心 O01 与 15 个受灾服务区的空间拓扑与分级生命线航路。

\begin{figure}[htbp]
  \centering
  \includegraphics[width=0.92\textwidth]{fig1_terrain_rescue_network.pdf}
  \caption{山区地形高程等高线与应急救援网络拓扑图}
  \label{fig:terrain}
\end{figure}

全流程技术路线涵盖三层递进架构流，如图~\ref{fig:roadmap}所示：
\begin{enumerate}
    \item \textbf{输入层（物理环境多维解析）}：读取 30m DEM 矩阵，建立高精地形高程场；解析 80 件货箱需求空间分布；构建爬升-巡航-下降三维飞行动力学能耗底座；
    \item \textbf{求解核（四阶段递进建模与协同求解引擎）}：
    \begin{itemize}
        \item \emph{阶段一（安全边界与精确组批）}：二分搜索最大物理载荷，位掩码 DP 求解 SPP 集合划分；
        \item \emph{阶段二（多点时空网络与流水线）}：Clarke-Wright 节约构造，改进 ALNS 启发式大邻域搜索，车电双时间轴防冲突排程；
        \item \emph{阶段三（通信射线追踪与中继优化）}：三维 LOS 遮挡检测，制高点优化选址，实体中继轮转调度与通信覆盖；
        \item \emph{阶段四（图论强绑定连通与独立配置）}：连通分支拓扑代数分析，分区独立运行峰值并发核算与缺口溯源；
    \end{itemize}
    \item \textbf{产出层（综合决策评估看板）}：输出标准 6-Sheet 官方交付数据集，开展全要素自动化质量审计，生成全套高精度学术图表资产。
\end{enumerate}

\begin{figure}[htbp]
  \centering
  \includegraphics[width=0.95\textwidth]{fig2_overall_methodology_roadmap.pdf}
  \caption{山区洪涝灾害无人机应急运输与通信协同优化技术路线与系统架构图}
  \label{fig:roadmap}
\end{figure}
